\documentclass[11pt]{article}
\usepackage[T1]{fontenc}
\usepackage{lmodern,amsmath,amssymb,amsthm,mathtools,booktabs,microtype}
\usepackage[margin=1in]{geometry}
\usepackage{xcolor}
\usepackage[colorlinks=true,linkcolor=blue!45!black,urlcolor=blue!45!black,citecolor=blue!45!black]{hyperref}
\usepackage{enumitem}
\newtheorem{theorem}{Theorem}[section]
\newtheorem{lemma}[theorem]{Lemma}
\newtheorem{corollary}[theorem]{Corollary}
\newtheorem{proposition}[theorem]{Proposition}
\theoremstyle{remark}\newtheorem{remark}[theorem]{Remark}
\newcommand{\R}{\mathbb R}\newcommand{\C}{\mathbb C}\newcommand{\supp}{\operatorname{supp}}\newcommand{\diam}{\operatorname{diam}}\newcommand{\res}{\operatorname{Res}}\newcommand{\Ree}{\operatorname{Re}}\newcommand{\Imm}{\operatorname{Im}}\newcommand{\dd}{\,\mathrm d}
\title{\textbf{Shortest robust positive filters for Dirichlet zeros}\\[4pt]\large Exact mirror-band solutions, local rigidity, and a spectral width differential}
\author{Research note for the ProveIt project}
\date{27 November 2026}
\begin{document}
\maketitle
\begin{abstract}
The ProveIt article asks for minimum-width positive averages annihilating several zero frequencies, including derivatives. We extend its consecutive-harmonic contour certificate to arbitrary mirror-symmetric spectral bands with multiplicities and complex exponents. We then prove a local classification beyond mirror symmetry: all sufficiently small independent perturbations of such a strictly positive configuration have a unique globally optimal measure up to translation, with exactly $n+1$ atoms whose locations and masses are real-analytic functions of the spectral parameters. Here $n$ is the total prescribed multiplicity in the upper half-plane. At the mirror configuration, the differential of the minimum width $D$ is
\[
 \dd D=-\frac{h^2}{\pi}\sum_j m_j\,\dd\gamma_j,
\]
where $h$ is the original grid spacing and $\gamma_j$ the imaginary parts. Real-part perturbations have zero first-order cost. In particular, the exact width for order $m$ cancellation at $0<\alpha<\gamma<3\alpha$ on a common vertical line is $4\pi m/(\alpha+\gamma)$; this is attained by a positive polynomial filter with $2m+1$ nodes. Dirichlet GRH makes this an optimal multiplicity-robust averaging theorem for explicit formulas. The extremal analysis is unconditional; the identification of actual zero contributions with critical-line frequencies uses Dirichlet GRH. We distinguish the new local theorem and differential from earlier two-mode input-shaping results.
\end{abstract}
\tableofcontents
\section{Source analysis and scope}
The inputs are the vendored readable Lean development \cite{vendor,main,statements,criterion} and the merged research article \cite{article,rec,zer}. The vendor README pins upstream commit \texttt{72fead102f663d5c7a407964a13dc0cf9d977207}, dated 25 November 2026, and states that it is an unmodified import of 818 files. Its top-level exported statements include \texttt{riemannHypothesis}, \texttt{LFunction\_ne\_zero\_of\_half\_lt\_re}, and \texttt{LFunction\_eq\_zero\_re}. The latter states that every zero of a Dirichlet $L$-function in $0<\Ree s<1$ lies on $\Ree s=1/2$; the half-plane formulation excludes zeros to its right (apart from the principal-character exception at the pole). These are statements about Dirichlet $L$-functions, not arbitrary Hecke, Artin or automorphic functions.

The upstream README reports successful kernel checks and independent exports with the standard three axioms \texttt{propext}, \texttt{Classical.choice}, and \texttt{Quot.sound}. The vendor README explicitly says the import itself was not rebuilt and is not part of ProveIt's root build. We inspected source declarations and the one-sided criterion, not a complete dependency replay; this note takes the exported Dirichlet GRH proposition as its analytic input. It does not certify its proof or claim independent community acceptance.

The source's final analytic step uses a positive smoothing of prime powers and a one-sided bound, with a twist chosen to eliminate contaminating zeros. In \texttt{Criterion.lean}, the kernel is $1/n-1/x$ for $n\le x$, and the Mellin residues encode orders of zeros. The Article develops a complementary transfer: positive dilation averages, whose Mellin multipliers can annihilate prescribed spectral contributions while preserving inequalities. Its Section \texttt{rec.tex}, subsection \texttt{rec:sec:dilation}, contains:
\begin{itemize}[leftmargin=*]
\item \texttt{rec:thm:fwidth}: sharp width for one frequency and arbitrary real parts and orders;
\item \texttt{rec:thm:hblock}: sharp width for consecutive harmonics with a common block of real parts and orders, proved using a keyhole-residue certificate;
\item \texttt{rec:thm:twofreq}: the simple-zero two-frequency classification; \texttt{rec:thm:cusp} gives its odd-resonance asymptotics;
\item \texttt{rec:thm:logfilter}: logarithmic support for large finite sets of ordinates, not exact minimum widths.
\end{itemize}
The open question at \texttt{rec:sec:open} asks about vanishing derivatives at two frequencies and the minimum widths for three noncommensurate frequencies. We give a complete answer for equal-order two-frequency constraints in the first ratio band, a more general exact theorem for mirror spectra, and an open-neighborhood solution for arbitrary independent perturbations, in particular for genuinely noncommensurate triples. We do not solve the global three-frequency problem or the higher odd-resonance problem.

The finite-spectrum optimization is not a standard consequence of GRH. The Article itself marks the underlying finite filter problem unconditional, and that is also true here. The arithmetic application is explicitly conditional on the exported theorem. The Article's README reports a merge of 131 research reports; it is not a verification of their results. Our proof below is self-contained, except for standard complex analysis, the real-analytic implicit function theorem, and the usual explicit formula for the application.

\section{The optimization problem}
For a compactly supported Borel probability measure $\mu$ on $\R$, put
\[
 F_\mu(s)=\int e^{su}\dd\mu(u),\qquad D(\mu)=\diam\supp\mu.
\]
For distinct $s_j=\beta_j+i\gamma_j$ with $\gamma_j>0$, and positive integers $m_j$, impose the jets
\begin{equation}\label{eq:jets}
 F_\mu^{(k)}(s_j)=\int u^ke^{s_ju}\dd\mu(u)=0\quad (0\le k<m_j).
\end{equation}
Write $n=\sum_jm_j$ and $D(\boldsymbol{s})$ for the infimum of diameters. Translation by $v$ multiplies $F$ by $e^{sv}$, preserving all prescribed orders. Conjugation supplies the lower-half-plane constraints automatically. Thus we may always normalize the left endpoint to $0$. No symmetry of an admissible measure is assumed. The set is nonempty, since one can convolve positive two-atom filters for each exponent and repeat them to the desired orders.

A \emph{certificate} for optimality on an interval $[0,L]$ is a real function $g$ in the real span of $\Ree(u^ke^{s_ju})$ and $\Imm(u^ke^{s_ju})$ that is nonnegative there. Every admissible measure supported there integrates $g$ to zero, hence is supported on its zero set. If the zero set is a finite grid with a unique feasible vector of positive endpoint masses, no shorter measure can be feasible. This observation turns the primal nonlinear optimization into a problem of constructing an exact dual certificate.

\section{An exact polynomial and contour theorem}
Fix $h>0$. In this section the coordinate $x=u/h$ is dimensionless. Set $a_j=hs_j$, $b_j=a_j-i\pi$, $r_j=e^{a_j}$ and $z_j=e^{b_j}=-r_j$. Assume:
\begin{enumerate}[label=(M\arabic*),leftmargin=*]
\item $0<\Imm a_j<2\pi$ for every $j$, so $-\pi<\Imm b_j<\pi$;
\item the multiset of $b_j$ with multiplicities $m_j$ is invariant under complex conjugation.
\end{enumerate}
Let
\begin{equation}\label{eq:PQ}
 P(z)=\prod_j(z-r_j)^{m_j},\quad Q(z)=\prod_j(z-z_j)^{m_j}=(-1)^nP(-z),\quad P(z)=\sum_{\ell=0}^np_\ell z^\ell.
\end{equation}
Both polynomials have real coefficients. Their roots are nonzero; $Q$ has no roots on the negative real axis. The $r_j$ are distinct since the imaginary parts of $a_j$ lie in an interval of length $2\pi$. Also $r_j\ne1$. In particular $P(1)\ne0$. In fact $P(1)>0$: a real root $r_j$ must be negative, and nonreal roots are paired.

\begin{theorem}[Exact mirror-band criterion]\label{thm:mirror}
Under (M1)--(M2), $D(\boldsymbol{s})\ge nh$. Equality holds if and only if all $p_\ell\ge0$. In that case the unique minimizer up to translation is
\begin{equation}\label{eq:optimal}
 \mu_* = \frac1{P(1)}\sum_{\ell=0}^np_\ell\,\delta_{\ell h}.
\end{equation}
If some $p_\ell<0$, the infimum is strictly larger than $nh$. In particular, if every $\Imm a_j\in(\pi/2,3\pi/2)$, all coefficients are strictly positive, and the optimal measure has $n+1$ atoms.
\end{theorem}
\begin{proof}
Define for real $x$
\[
 E(x)=\sum_j\res_{z=z_j}\frac{z^{x-1}}{Q(z)},
\]
using the principal branch of $\log z$ on the slit plane. Expanding the residue at a root of order $m_j$ expresses it as $e^{b_jx}$ times a polynomial of degree at most $m_j-1$. Conjugation symmetry makes $E(x)$ real. Therefore
\begin{equation}\label{eq:g}
 g(x)=(-1)^{n+1}\sin(\pi x)E(x)
\end{equation}
belongs to the real span of $\Ree(x^k e^{a_jx})$ and $\Imm(x^k e^{a_jx})$: since $E$ is real, $\sin(\pi x)E(x)=\Imm(e^{i\pi x}E(x))$.

For $0<x<n$ integrate on a keyhole about the negative real axis, using radii $\epsilon$ and $R$ with $\epsilon\to0$, $R\to\infty$. The small-circle integral is $O(\epsilon^x)$ and the large-circle integral is $O(R^{x-n})$. The upper and lower values of $z^{x-1}$ on the cut are $-t^{x-1}e^{\pm i\pi x}$. Moreover $Q(-t)$ has sign $(-1)^n$ for $t>0$, because real roots of $Q$ are positive and conjugate pairs give positive products. The residue theorem, with the upper cut traversed from $-R$ to $0$, gives
\begin{equation}\label{eq:keyhole}
 E(x)=\frac{(-1)^{n+1}\sin(\pi x)}{\pi}\int_0^\infty\frac{t^{x-1}}{|Q(-t)|}\dd t,
 \qquad
 g(x)=\frac{\sin^2(\pi x)}{\pi}\int_0^\infty\frac{t^{x-1}}{|Q(-t)|}\dd t.
\end{equation}
The integral is finite and strictly positive. Hence the only interior zeros of $g$ on $[0,n]$ are the integers, and each is exactly double. Since $E$ is an entire exponential polynomial in $x$, $g(0)=g(n)=0$. In addition,
\begin{equation}\label{eq:slopes}
 g'(0)=\pi/|Q(0)|,\qquad g'(n)=-\pi.
\end{equation}
Indeed near $x=0$ the integral in \eqref{eq:keyhole} is $1/(|Q(0)|x)+O(1)$, by separating $(0,1)$ and using $|Q(-t)|^{-1}=|Q(0)|^{-1}+O(t)$. Near $x=n$ it is $1/(n-x)+O(1)$ by $|Q(-t)|^{-1}=t^{-n}(1+O(1/t))$ for large $t$. This also proves simple endpoint zeros.

Translate any candidate of width at most $nh$ into $[0,nh]$ and scale it into $[0,n]$. Its $g$ integral vanishes by \eqref{eq:jets}; thus its masses are supported on $\{0,1,\ldots,n\}$. Its weight polynomial $A(z)=\sum_{\ell=0}^n w_\ell z^\ell$ must have zeros at every $r_j$ of order at least $m_j$, since its transform is $A(e^{hs})$ and the exponential has nonzero derivative there. Consequently $A=P/P(1)$. The constant term and leading coefficient of $P$ are strictly positive, so both endpoints are required. A feasible polynomial exists precisely when all its coefficients are nonnegative, and then it satisfies every constraint. This proves the bound, equality and uniqueness. If negative coefficients prevent equality but diameters tend to $nh$, translate to left endpoint $0$ and take a weak limit on a common compact interval. The moment conditions pass to the limit by continuity and the limit is supported on $[0,nh]$, a contradiction. Finally in the open band $\Ree r_j<0$, real-root factors have strictly positive coefficients, and pair factors $z^2-2\Ree(r_j)z+|r_j|^2$ do as well. Their product has every coefficient strictly positive.
\end{proof}
\begin{remark}
The nonnegative-coefficient criterion can apply beyond the open band. At its boundary, some coefficients may vanish; the exact theorem still applies, but the perturbation theorem below requires strictly positive masses. The result includes the consecutive-harmonic theorem of the Article and allows noncommensurate pairs of reflected frequencies, varying multiplicities from pair to pair, and varying real parts. The keyhole step is adapted from the Article, not a claim of a newly discovered contour method.
\end{remark}

\begin{corollary}[Equal-order pair, first band]\label{cor:pair}
Let $0<\alpha<\gamma<3\alpha$, $\beta\in\R$ and $m\ge1$. Require order $m$ at both $\beta+i\alpha$ and $\beta+i\gamma$. Then
\[
 D_m=\frac{4\pi m}{\alpha+\gamma},\qquad h=\frac{2\pi}{\alpha+\gamma},\qquad
 \sum_{\ell=0}^{2m}w_\ell z^\ell=\frac{(z^2-2e^{h\beta}\cos(h\alpha)z+e^{2h\beta})^m}{(1-2e^{h\beta}\cos(h\alpha)+e^{2h\beta})^m}.
\]
The unique normalized minimizer is $\sum_\ell w_\ell\delta_{\ell h}$. Its width is a factor $4\alpha\gamma/(\alpha+\gamma)^2$ of the width $m\pi(1/\alpha+1/\gamma)$ of two separately designed filters. At $\gamma=3\alpha$ the same formula with zero odd coefficients gives width $m\pi/\alpha$. The endpoint is not a strictly positive configuration.
\end{corollary}
\begin{proof}
The angles $h\alpha$ and $h\gamma$ sum to $2\pi$; the first lies in $(\pi/2,\pi)$, hence the mirror conditions and positivity hold. The boundary quadratic has no linear term, but its other coefficients are positive. Compare the two explicit widths.
\end{proof}

\section{Local solution for nonsymmetric spectra}
The principal extension is the following theorem. The parameters $s_j$ are independent complex numbers: the local result does not require a mirror pairing or any exact rational relations after perturbation. For fixed $m_j$, the neighborhood is in $\R^{2J}$, where $J$ is the number of distinct upper-half-plane exponents. It may depend on the base configuration.

\begin{theorem}[Analytic continuation of the global optimizer]\label{thm:local}
Suppose the base $\boldsymbol{s}^0$ satisfies (M1)--(M2) and all coefficients in \eqref{eq:PQ} are strictly positive. There is an open neighborhood $U$ of $\boldsymbol{s}^0$ in which every minimum-width problem \eqref{eq:jets} has a unique minimizer up to translation. It has the form
\begin{equation}\label{eq:localmeasure}
 \mu_{\boldsymbol{s}}=\sum_{\ell=0}^n w_\ell(\boldsymbol{s})\,\delta_{h x_\ell(\boldsymbol{s})},\quad
 0=x_0<x_1<\cdots<x_n,\quad w_\ell>0,\quad \sum w_\ell=1,
\end{equation}
where the masses, nodes and $D(\boldsymbol{s})=h x_n(\boldsymbol{s})$ are real-analytic on $U$. A real-analytic family of dual certificates is positive between the nodes, has double interior zeros and simple endpoint zeros. In particular the optimizer is globally minimal among all compactly supported probabilities, not only among atomic measures with $n+1$ atoms.
\end{theorem}

\subsection{The nonsingular moment Jacobian}
We use the following elementary fact about sampled exponential polynomials.
\begin{lemma}[Confluent sampling]\label{lem:sample}
Let distinct nonzero $z_j$ have multiplicities $m_j$, with total $n$. The sequences $k^r z_j^k$ for $0\le r<m_j$ form a basis of solutions of the recurrence $\prod_j(S-z_j)^{m_j}y=0$, where $S$ is the forward shift. Evaluation at any $n$ consecutive integers is an isomorphism. If no $z_j=1$, adjoining the constant sequence gives the analogous assertion for $n+1$ samples.
\end{lemma}
\begin{proof}
Each sequence is annihilated by the recurrence, and they are independent: the generating functions $\sum_{k\ge0}k^r z_j^k t^k$ have poles at $t=1/z_j$ of exact order $r+1$, with a nonzero leading coefficient; inspecting highest poles successively forces any linear dependence to vanish. The order-$n$ monic recurrence has nonzero constant coefficient, so $n$ successive values determine the solution both forwards and backwards. The constant sequence is the extra root $1$.
\end{proof}

Fix $h$ at its base value. Let $V$ be the real span of $x^k e^{b_jx}$ (real combinations in the sense of taking the real-valued part of the complex span). Its real dimension is $n$ by conjugation closure and independence of exponential polynomials. Let $W$ be the real span of the $2n$ jet functions $\Ree(x^ke^{a_jx})$, $\Imm(x^ke^{a_jx})$ at the base. Then
\begin{equation}\label{eq:decomposition}
 W=\cos(\pi x)V\oplus\sin(\pi x)V.
\end{equation}
Indeed the real and imaginary parts of the shifted exponential polynomial lie in $V$, so the two spaces sum to $W$; the $2n$ complex exponents $a_j,\bar a_j$ are distinct between the two halves and independent with their polynomial factors. This gives the dimension and directness. For clarity, independence of continuous exponential polynomials follows, for example, by applying products of the differential operators $(\frac{\dd}{\dd x}-a_j)^{m_j}$ to isolate one exponent at a time. On integer arguments $\ell$, a function $f=c+\cos(\pi x)A(x)+\sin(\pi x)B(x)$ in $\R+W$ has values $c+(-1)^\ell A(\ell)$.

Consider the square real system with unknowns $w_0,\ldots,w_n,x_1,\ldots,x_n$ and $x_0=0$:
\begin{equation}\label{eq:system}
 \sum_{\ell=0}^n w_\ell=1,\qquad
 \sum_{\ell=0}^n w_\ell x_\ell^k e^{hs_jx_\ell}=0\quad(0\le k<m_j),
\end{equation}
with complex equations split into real and imaginary parts. Replacing physical powers by $x^k$ changes only nonzero scale factors. The Jacobian at the base has columns given by evaluation of the $2n+1$ functions spanning $\R+W$ at all nodes and derivative evaluation at $1,\ldots,n$ times $w_\ell^0>0$.

A left null vector corresponds to $f=c+\cos(\pi x)A+\sin(\pi x)B$ vanishing at $0,\ldots,n$ and with $f'(\ell)=0$ for $1\le\ell\le n$. The values form an exponential sequence with roots $1,-z_j$, which are distinct since $r_j=-z_j\ne1$. Lemma~\ref{lem:sample} gives $c=0$ and $A=0$. Then $f'(\ell)=\pi(-1)^\ell B(\ell)$, so $B$ vanishes at $n$ consecutive integers; the same lemma gives $B=0$. Thus the Jacobian is invertible. The real-analytic implicit function theorem supplies unique local analytic masses and nodes, which remain strictly positive and ordered after shrinking the neighborhood. These give a feasible measure, but optimality still requires the next step.

\subsection{Persistence of the positive dual certificate}
For perturbed exponents define $W_{\boldsymbol{s}}$ by the same real and imaginary jet functions, with $a_j=hs_j$. Seek $g_{\boldsymbol{s}}\in W_{\boldsymbol{s}}$ such that
\begin{equation}\label{eq:dualconds}
 g_{\boldsymbol{s}}(x_\ell)=0\ (0\le\ell\le n),\qquad
 g'_{\boldsymbol{s}}(x_\ell)=0\ (1\le\ell<n).
\end{equation}
This is a homogeneous $2n\times 2n$ linear system. At the base, a function in its kernel has $A=0$ in \eqref{eq:decomposition} by confluent sampling, and $B(1)=\cdots=B(n-1)=0$. These conditions on $B$ have rank $n-1$ (extend to $n$ consecutive samples). Thus the kernel is one-dimensional, spanned by \eqref{eq:g}. For nearby spectra the rank is at least $2n-1$ by a nonvanishing minor. It is at most $2n-1$, since the feasible positive weights from \eqref{eq:system} give the nontrivial row relation $\sum w_\ell g(x_\ell)=0$ for all $g\in W_{\boldsymbol{s}}$. In particular one can discard the row at $x_0$ at the base and nearby, since $w_0>0$ and the rank defect is one. Normalize by a fixed coefficient functional that takes a nonzero value on the base certificate. The resulting linear solve gives an analytic certificate with $g_{\boldsymbol{s}^0}=g$.

At the base, \eqref{eq:keyhole} and \eqref{eq:slopes} show strict positivity away from the nodes, positive second derivative at interior nodes, positive derivative at the left endpoint and negative derivative at the right. These properties persist: on fixed small neighborhoods of the original interior nodes the second derivative stays uniformly positive and the imposed double-zero conditions give $g_{\boldsymbol{s}}(x)>0$ except at the shifted nodes. The endpoint derivatives remain strictly signed on their small neighborhoods; use the imposed endpoint zeros and the mean value theorem. On the remaining compact pieces uniform convergence preserves a strict positive lower bound. Thus $g_{\boldsymbol{s}}\ge0$ on $[0,x_n]$ with precisely the displayed nodes as zeros.

Any admissible measure of diameter at most $h x_n$ can be translated and scaled into this interval. Its certificate integral is zero, so it is supported on the same nodes. The evaluation columns of the moment Jacobian remain linearly independent by openness of its nonsingularity; hence the mass vector satisfying the moments and normalization is unique. In particular both endpoints must be occupied. The solution is globally optimal and unique after fixing the left endpoint. This completes the proof of Theorem~\ref{thm:local}.

\section{A universal first-order spectral width law}
\begin{theorem}[Differential at every strictly positive mirror configuration]\label{thm:differential}
Under the hypotheses of Theorem~\ref{thm:local}, write $s_j=s_j^0+\Delta\beta_j+i\Delta\gamma_j$. Then
\begin{equation}\label{eq:diff}
 D(\boldsymbol{s})=nh-\frac{h^2}{\pi}\sum_jm_j\Delta\gamma_j+O_{\boldsymbol{s}^0,\boldsymbol m}(\|\Delta\boldsymbol{s}\|^2).
\end{equation}
All real-part derivatives vanish, even when the different real parts vary independently. The leading change in width depends only on the multiplicity-weighted sum of imaginary-part perturbations, not on the original root locations within the band.
\end{theorem}
\begin{proof}
Work in dimensionless coordinates. At the base, expand the residue polynomial as
\[
 E(x)=\sum_j e^{b_jx}\sum_{k=0}^{m_j-1}c_{jk}x^k,
 \qquad c_{j,m_j-1}=\frac{1}{(m_j-1)!z_j^{m_j}\prod_{v\ne j}(z_j-z_v)^{m_v}}.
\]
The leading coefficient follows by differentiating $z^{x-1}$ exactly $m_j-1$ times in the residue formula. Freeze all $c_{jk}$ at the base. Along a smooth parameter path $a_j(\epsilon)$ with $a_j(0)=a_j$, the real function
\[
 G_\epsilon(x)=(-1)^{n+1}\Imm\sum_{j,k}c_{jk}x^k e^{a_j(\epsilon)x}
\]
is a linear combination of the current moment functions; $G_0=g$ by \eqref{eq:g}. Thus $\sum_\ell w_\ell(\epsilon)G_\epsilon(x_\ell(\epsilon))=0$. Differentiate at zero. The terms containing $\dot w_\ell$ vanish since $g(\ell)=0$, the interior-node terms vanish since $g'(\ell)=0$, and the left endpoint is fixed. By \eqref{eq:slopes},
\begin{equation}\label{eq:derivative1}
 0=\sum_\ell w_\ell^0\dot G(\ell)-\pi w_n^0\dot x_n.
\end{equation}
At $r_j=e^{a_j}=-z_j$, the moment identity gives
\[
 \sum_\ell w_\ell^0 \ell^p r_j^\ell
 =\frac{(z\partial_z)^pP(r_j)}{P(1)}.
\]
For $p<m_j$ this is zero, and for $p=m_j$ its numerator is $r_j^{m_j}m_j!\prod_{v\ne j}(r_j-r_v)^{m_v}$. The derivative $\dot G$ raises the degree $k$ by one, so only $k=m_j-1$ survives. Multiplication by the leading residue coefficient yields
\[
 c_{j,m_j-1}r_j^{m_j}m_j!\prod_{v\ne j}(r_j-r_v)^{m_v}=(-1)^n m_j.
\]
Therefore $\sum_\ell w_\ell^0\dot G(\ell)=-\Imm\sum_jm_j\dot a_j/P(1)$. As $w_n^0=1/P(1)$, \eqref{eq:derivative1} becomes $\dot x_n=-\Imm\sum_jm_j\dot a_j/\pi$. Finally $D=h x_n$ with fixed $h$, so $\dot D=-h^2\sum_jm_j\dot\gamma_j/\pi$. Real analyticity supplies the stated quadratic remainder.
\end{proof}

For example, rescaling every frequency by $1+\epsilon$ and every real part similarly gives $D/(1+\epsilon)$ exactly. At the base $\sum_jm_j\gamma_j=n\pi/h$ by mirror symmetry, and \eqref{eq:diff} reproduces the derivative $-nh$. More generally, keeping the weighted mean imaginary part fixed makes the first-order width cost vanish; one can independently displace all real parts and destroy commensurability without paying any linear cost. The theorem does not claim that the second-order term always has a particular sign.

\section{Stability and spectral applications}
\begin{proposition}[Local sharp error stability]\label{prop:stability}
Shrink $U$ in Theorem~\ref{thm:local} if necessary. For $\boldsymbol{s}\in U$ and a probability $\nu$ supported on $[0,D(\boldsymbol{s})]$, let $\varepsilon$ be the maximum absolute value of its prescribed jets. There is a uniform constant $C$ such that for $0\le\varepsilon\le1$
\[
 W_1(\nu,\mu_{\boldsymbol{s}})\le C\varepsilon^{1/2}\quad(n\ge2),
\]
where $W_1$ is the transport distance with cost $|u-v|$. The exponent is sharp at every such configuration. If $n=1$, the sharp exponent is $1$.
\end{proposition}
\begin{proof}
The normalized certificates vary continuously, and the endpoint and interior derivative bounds from the local proof can be chosen uniform in a smaller neighborhood. Thus $g_{\boldsymbol{s}}(u/h)\ge c\operatorname{dist}(u,G)^2$ for $n\ge2$ (and $\ge c\operatorname{dist}(u,G)$ for $n=1$), where $G$ is the optimal node set. Its integral is bounded by $C\varepsilon$ using the bounded coefficients of the certificate in the jet basis. Move every mass of $\nu$ to a nearest node, which costs $O(\sqrt\varepsilon)$ by Cauchy--Schwarz, or $O(\varepsilon)$ for $n=1$. On the common bounded interval the moment functions are uniformly Lipschitz, so the rounded masses have moment residuals of this order. The evaluation matrix augmented by the mass row has full column rank, uniformly with a bounded left inverse after shrinking $U$. Thus the rounded vector is within this order of the unique weights, and redistributing it costs at most the interval diameter times its total variation discrepancy. To prove sharpness for $n\ge2$, split one positive interior mass equally at distance $\eta$ to either side, holding the other masses fixed. Every jet changes by $O(\eta^2)$, and the transport cost is exactly the moved mass times $\eta$ for sufficiently small $\eta$. For $n=1$, move one endpoint mass inward by $\eta$; jets change by $O(\eta)$ and transport is of order $\eta$.
\end{proof}

\begin{corollary}[Optimal robust explicit-formula cancellation; uses Dirichlet GRH]\label{cor:prime}
Let $L(s,\chi)$ be a Dirichlet $L$-function, and choose two distinct positive nontrivial zero ordinates $\alpha,\gamma$ from it or from different characters, with $\gamma<3\alpha$. For a smooth compactly supported weight $W$ on $(0,\infty)$, consider the smoothed prime-power sum $S_{\chi,W}(X)=\sum_n\Lambda(n)\chi(n)W(n/X)$. Under the exported Dirichlet GRH, the two zeros lie on $\Ree s=1/2$. Among positive probability averages of $S_{\chi,W}(Xe^u)$ whose Mellin multipliers vanish to order $m$ at both zeros, the exact least multiplicative range is
\[
 \frac{X_{\max}}{X_{\min}}=\exp\left(\frac{4\pi m}{\alpha+\gamma}\right),
\]
with weights in Corollary~\ref{cor:pair} for $\beta=1/2$. This works simultaneously for any collection of such sums, independent of their Mellin weights. Conjugate zeros are canceled as well when they occur in the formulas (possibly in the conjugate character's function).
\end{corollary}
\begin{proof}
The smoothed explicit formula gives a contribution $-m_\rho \widetilde W(\rho)X^\rho$ from a zero of multiplicity $m_\rho$; this follows from Mellin inversion of $-L'/L$ and residues. The dilation average replaces it by the same expression times $F_\mu(\rho)$. All contributions associated with a target zero, including formal derivatives up to order $m-1$, vanish by the multiplier conditions. The finite-width assertion is Corollary~\ref{cor:pair}. This optimality concerns the prescribed multiplier constraints, not special weights already vanishing at target zeros. If a principal term occurs, it is multiplied by $F_\mu(1)>0$ and can be accounted for or rescaled. The remaining zeros, trivial terms and pole terms are not claimed to vanish.
\end{proof}

Order $m$ also guarantees robustness if the zero is measured with a small complex error $\eta$: for a normalized filter supported on $[0,D]$,
\begin{equation}\label{eq:robust}
 |F_\mu(s_j+\eta)|\le e^{\max(\Ree s_j,0)D}e^{|\eta|D}\frac{(|\eta|D)^{m_j}}{m_j!}.
\end{equation}
Indeed expand $e^{\eta u}$ through degree $m_j-1$ and bound the Taylor remainder uniformly. This estimate addresses robustness to ordinate errors (or real-part errors) separately from the weight perturbation stability in Proposition~\ref{prop:stability}. The cancellation of all polynomial factors in $\log X$ up to degree $m_j-1$ follows by binomial expansion of $(\log X+u)^k$.

\section{Examples and reproducibility}
\subsection{The first two zeta ordinates}
The illustrative high-precision values $\alpha\approx14.134725141735$, $\gamma\approx21.022039638772$ satisfy the first-band inequality. For $m=2$ the minimum width is $0.7148763939353$ and the scale ratio is $2.043934024003$; separate filters would require width $0.7434068409396$. Direct Mellin weights at $\beta=1/2$ are as follows. Numbers are rounded; use the exact expressions, not the rounded table, for cancellation.
\begin{center}
\begin{tabular}{rr}
\toprule
Shift $u$ & Probability mass\\\midrule
0 & 0.090192484288\\
0.178719098484 & 0.269392994857\\
0.357438196968 & 0.352023599548\\
0.536157295451 & 0.225304351702\\
0.714876393935 & 0.063086569606\\\bottomrule
\end{tabular}
\end{center}
The same formula applies to any order $m$, without solving an optimization problem or making any assumption on linear independence of ordinates. It does not prove simplicity or separation of zeros from different $L$-functions.

\subsection{A truly nonsymmetric three-frequency perturbation}
In dimensionless units $h=1$, take the base exponents $i(2\pi/3,\pi,4\pi/3)$, all of order one. Then $P(z)=(z+1)(z^2+z+1)$, and the optimizer has masses $(1,2,2,1)/6$ at $0,1,2,3$. Perturb to
\[
 a_1=2\pi i/3+\epsilon/3,\qquad a_2=i(\pi+\epsilon),\qquad a_3=4\pi i/3.
\]
For all sufficiently small real $\epsilon$, the unique shortest filter is described by Theorem~\ref{thm:local}, and its width is $3-\epsilon/\pi+O(\epsilon^2)$. In particular if $\epsilon$ is nonzero rational, $a_2$'s frequency is incommensurate with the first (its ratio is $3/2+3\epsilon/(2\pi)$, irrational). The same width differential holds for the all-imaginary perturbation obtained by omitting $\epsilon/3$. This example is a concrete open-neighborhood theorem rather than a statement restricted to commensurate frequencies. Its proof does not provide an explicit radius of the neighborhood.

\begin{center}
\begin{tabular}{rrrr}
\toprule
$\epsilon$ & Computed width & $3-\epsilon/\pi$ & Difference$/\epsilon^2$\\\midrule
0.001 & 2.999681733 & 2.999681690 & 0.04293\\
0.0001 & 2.999968169 & 2.999968169 & 0.04294\\
-0.001 & 3.000318353 & 3.000318310 & 0.04295\\\bottomrule
\end{tabular}
\end{center}
The numerical continuation solves all real and imaginary constraints and evaluates the dual at a dense sample. These are diagnostic checks only: positivity and optimality in an unspecified sufficiently small neighborhood follow analytically, not from the sample. The supplied \texttt{verify.py} generates \texttt{results.json} with 65-digit arithmetic (values printed to 32 significant digits), including weights, nodes, and residuals. It uses \texttt{mpmath.zetazero}; this is not a certified zero isolation routine. No computation is required for the mathematical theorems.

\section{Prior art, significance, and further questions}
A focused literature search found robust multi-mode input shaping, where the same positive-impulse optimization arises for vibration cancellation \cite{lau,pao,command}. In particular, the accessible text of \emph{Command Generation for Flexible Systems}, Section 3.4.2, credits Rappole with three- and five-impulse closed forms for two undamped modes, and discusses time-optimality in the range of frequency ratios up to three. Thus we do not claim priority for the order-one or order-two pair formulas in isolation. The Article already proves the one-frequency and consecutive-harmonic contour bounds. Our candidate contribution is the extension to arbitrary mirror spectra with variable multiplicities and the local analytic classification and explicit differential for independently perturbed complex exponents, established by a confluent-sampling nonsingularity argument and persistence of a positive dual. No matching theorem for that full combination was identified in this limited search; this is not an exhaustive priority determination.

The result addresses part of the Article's explicitly posed shortest-positive-average problem. It is useful because exact cancellation alone does not guarantee minimum width, and because numerical optimizers that restrict the number of atoms or impose symmetry need not establish global optimality. Here the certificate establishes both a unique support and sharp width. The local theorem moreover permits nonsymmetric and noncommensurate input data while guaranteeing the same atom count. The differential gives a directly computable sensitivity law independent of fine spacings within the base mirror band. Under Dirichlet GRH the theorem applies to real spectral frequencies of prime errors, with arbitrary prescribed vanishing order and no linear-independence hypothesis. It does not resolve any major named conjecture outside the stated optimization problem.

Further research questions:
\begin{enumerate}[leftmargin=*]
\item Make the neighborhood in Theorem~\ref{thm:local} effective in terms of root separation, minimum weight, and certificate curvature; use interval arithmetic to certify particular nonmirror zero triples.
\item Determine the Hessian of $D$ at mirror configurations and its sign on weighted-mean-preserving frequency perturbations. The first-order law is exact but does not settle these directions.
\item Determine the global phase diagram when the dual certificate acquires additional zeros or a mass tends to zero. In particular, solve equal-order pairs above ratio $3$, and general triples outside the local neighborhood.
\item Optimize joint cancellation of all low zero ordinates with nonnegative weights, including the leading constant in the Article's logarithmic-support theorem; the fixed-band theorem does not settle this asymptotic problem.
\item Formalize the finite sampling and residue certificate in Lean, then combine with the exported Dirichlet GRH theorem and a formal explicit formula. The current note contains a classical proof and numerical checks, not a Lean-checked extension.
\end{enumerate}

\begin{thebibliography}{99}
\bibitem{vendor} ProveIt, \emph{Vendored grh-lean (readable edition)}, README, blob \texttt{a04f66a62af93fc0757f2aa573f3b053fde51bfa}. \href{https://github.com/VladimirReshetnikov/ProveIt/tree/main/Analysis/GeneralizedRiemannHypothesis/vendor/grh-lean}{Repository directory}.
\bibitem{main} L. Alp\"oge with Claude; lsernau/grh-lean, \emph{src/Main.lean}, vendored blob \texttt{1adf4187a938510363cdb0dde0a6d628ff39c600}. \href{https://github.com/VladimirReshetnikov/ProveIt/blob/main/Analysis/GeneralizedRiemannHypothesis/vendor/grh-lean/grh/src/Main.lean}{Source}.
\bibitem{statements} Same, \emph{src/Statements.lean}, vendored blob \texttt{1cd1bce93344065936b2b392b23306254e0785ef}. \href{https://github.com/VladimirReshetnikov/ProveIt/blob/main/Analysis/GeneralizedRiemannHypothesis/vendor/grh-lean/grh/src/Statements.lean}{Source}.
\bibitem{criterion} Same, \emph{src/New/Reduction/Zeros/Criterion.lean}, vendored blob \texttt{5c06362a50e044d491423fb3d830bae6a50c0075}. \href{https://github.com/VladimirReshetnikov/ProveIt/blob/main/Analysis/GeneralizedRiemannHypothesis/vendor/grh-lean/grh/src/New/Reduction/Zeros/Criterion.lean}{Source}.
\bibitem{article} ProveIt, \emph{Beyond the Riemann hypothesis for Dirichlet L-functions: consequences, transfers and limits}, Article README, blob \texttt{7f3286088c94a8df3d8b9846112ea83dbbbbfa25}. \href{https://github.com/VladimirReshetnikov/ProveIt/tree/main/Analysis/GeneralizedRiemannHypothesis/Article}{Article directory}.
\bibitem{rec} Same, \emph{sections/rec.tex}, blob \texttt{da56d6ad5a1a699e335dcfd327bcdc9492c64017}. \href{https://github.com/VladimirReshetnikov/ProveIt/blob/main/Analysis/GeneralizedRiemannHypothesis/Article/sections/rec.tex}{Source}.
\bibitem{zer} Same, \emph{sections/zer.tex}, blob \texttt{1f1e31367ea4fdd259786e3997d46ecd279d5cbc}. \href{https://github.com/VladimirReshetnikov/ProveIt/blob/main/Analysis/GeneralizedRiemannHypothesis/Article/sections/zer.tex}{Source}.
\bibitem{lau} M. A. Lau and L. Y. Pao, Input shaping and time-optimal control of flexible structures, \emph{Automatica} 39 (2003), 893--900. \href{https://doi.org/10.1016/S0005-1098(03)00024-4}{doi:10.1016/S0005-1098(03)00024-4}.
\bibitem{pao} L. Y. Pao and W. E. Singhose, Robust minimum time control of flexible structures, \emph{Automatica} 34 (1998), 229--236. \href{https://doi.org/10.1016/S0005-1098(97)00178-7}{doi:10.1016/S0005-1098(97)00178-7}.
\bibitem{command} \emph{Command Generation for Flexible Systems}, Sec. 3.4.2, accessible text via CiteSeerX; source of the attribution to Rappole. \href{https://citeseerx.ist.psu.edu/document?doi=7aae7b1b0fce193f503c9cc523af9972fcf16e68\&repid=rep1\&type=pdf}{Document}.
\end{thebibliography}
\end{document}
